\chapter*{Exercises}
\phantomsection
\addcontentsline{toc}{chapter}{Exercises}

\begin{enumerate}
\def\labelenumi{\arabic{enumi})}
\item
  With the hypotheses of No.~1, show that the set of bounded functions on X such that \(M(|f|)\) is finite is a subalgebra A of \(\mathbf{R}^X\) , and that the set of bounded functions on X such that \(M(|f|) = 0\) is an ideal in A. If, moreover, \(M(1)\) is finite, show that the mapping \(f \mapsto M(f)\) is continuous when A is equipped with the topology of uniform convergence in X.
\item
  Let X be the interval \([0, +\infty[\) of R, S the convex cone formed by the functions defined on X such that \(0 \leqslant f(x) \leqslant kx\) on X (for a finite number k \textgreater{} 0 depending on f). Set \(M(f) = 0\) for \(f \in S\) , and \(M(f) = +\infty\) for every positive function f defined on X and not belonging to S. Show that M satisfies the conditions of No.~1, and that \(M(x) = 0\) , and \(M(x^{r}) = +\infty\) for every number r \textgreater{} 0 not equal to 1.
\item
  Give an example where \(X\) is a set with two elements, \(N sub p(\mathbf{x})\) is finite for every \(p > 0\) and every \(\mathbf{x} \in \mathbf{R}^2\) , but where there exist values of \(p\) such that the mapping \(p \mapsto N sub p(\mathbf{x})\) is not differentiable at these points.
\item
  Deduce the inequality (6) from the inequality of the geometric mean
\end{enumerate}

\[
z _ {1} ^ {\alpha_ {1}} z _ {2} ^ {\alpha_ {2}} \dots z _ {n} ^ {\alpha_ {n}} \leqslant \alpha_ {1} z _ {1} + \dots + \alpha_ {n} z _ {n} \quad (\text {where} \sum_ {i = 1} ^ {n} \alpha_ {i} = 1)
\]

(FRV, III, §1, No.~1, Prop. 2). (Reduce to the case that \(M(f_i) = 1\) for \(1 \leqslant i \leqslant n\) .)

\begin{enumerate}
\def\labelenumi{\arabic{enumi})}
\setcounter{enumi}{4}
\tightlist
\item
  Let \(\alpha\) be a real number \(>1\) and let \(\beta = 1 - \alpha < 0\) . Let \(g\) be a finite function, defined on X, such that \(g(x) > 0\) for all \(x \in X\) and such that \(M(g) > 0\) ; show that for every finite function \(f \geqslant 0\) defined on X such that \(M(f)\) is finite,
\end{enumerate}

\[
M (f ^ {\alpha} g ^ {\beta}) \geqslant (M (f)) ^ {\alpha} (M (g)) ^ {\beta}
\]

(apply Hölder's inequality suitably).

\begin{enumerate}
\def\labelenumi{\arabic{enumi})}
\setcounter{enumi}{5}
\tightlist
\item
  Deduce Minkowski's inequality from Hölder's inequality (find an upper bound for \(M(f(f + g)^{p-1})\) with the help of Hölder's inequality). If one assumes that \(M(f + g) = M(f) + M(g)\) for every pair of functions \(f, g\) defined and \(\geqslant 0\) on \(X\) , deduce similarly from Exer. 5 the inequality
\end{enumerate}

\[
\left(M \big ((f + g) ^ {p} \big)\right) ^ {1 / p} \geqslant \left(M (f ^ {p})\right) ^ {1 / p} + \left(M (g ^ {p})\right) ^ {1 / p}
\]

in the following cases:

\begin{enumerate}
\def\labelenumi{\alph{enumi})}
\tightlist
\item
  \(0 < p < 1\) , \(f\) and \(g\) finite functions \(\geqslant 0\) defined on X, such that \(f(x) + g(x) > 0\) for all \(x \in X\) and such that \(M(f^p)\) and \(M(g^p)\) are finite; b) \(p < 0\) , \(f\) and \(g\) finite functions defined on X, such that \(f(x) > 0\) and \(g(x) > 0\) for all \(x \in X\) , \(M(f^p)\) and \(M(g^p)\) are finite, and \(M((f + g)^p) > 0\) .
\end{enumerate}

